% ---
% titre: Tableau de critères de divisibilité
% theme: NO
% chapitre: Nombres naturels et décimaux
% sous_chapitre: Critères de divisibilité
% niveau: [9e]
% type: EVAL
% tags: [cc-criteres-divisibilite, p-question-TS]
% difficulte: 1
% corrige: true
% ---

\question[4]
 Complète le tableau suivant en notant une croix aux endroits qui conviennent.
 
\vspace{10pt}\def\arraystretch{2}
\begin{tabularx}{\linewidth}{|r*{8}{|>{\centering\arraybackslash}X}|}
\hline
\rowcolor{gray!20}\textbf{Divisible par \ldots} & \textbf{2} & \textbf{3} & \textbf{4} & \textbf{5} & \textbf{6} & \textbf{9} & \textbf{10} & \textbf{25}\\
\hline
156 & \textcolor{\colSolution} X & \textcolor{\colSolution} X & \textcolor{\colSolution} X & & \textcolor{\colSolution} X & & & \\
\hline
7\,587 & & \textcolor{\colSolution} X & & & & \textcolor{\colSolution} X & & \\
\hline
2\,100 & \textcolor{\colSolution} X & \textcolor{\colSolution} X & \textcolor{\colSolution} X & \textcolor{\colSolution} X & \textcolor{\colSolution} X & & \textcolor{\colSolution} X & \textcolor{\colSolution} X \\
\hline
716\,808 & \textcolor{\colSolution} X & \textcolor{\colSolution} X & \textcolor{\colSolution} X & & \textcolor{\colSolution} X & & & \\
\hline
\end{tabularx}
\def\arraystretch{1}

\begin{solution}[0pt]
\textbf{0.5pt} par colonne correcte
\end{solution}

